% Generated draft: requires mathematical and source review.
\documentclass[11pt]{article}
\usepackage[margin=25mm]{geometry}
\usepackage{fontspec}
\setmainfont{Latin Modern Roman}
\usepackage{amsmath,amssymb,mathtools,unicode-math}
\setmathfont{Latin Modern Math}
\usepackage{enumitem,needspace,xurl,hyperref}
\setlength{\parindent}{0pt}
\setlength{\parskip}{5pt}
\setlength{\emergencystretch}{3em}
\newcommand{\N}{\mathbb{N}}
\newcommand{\Z}{\mathbb{Z}}
\newcommand{\Q}{\mathbb{Q}}
\newcommand{\R}{\mathbb{R}}
\newcommand{\C}{\mathbb{C}}
\newcommand{\F}{\mathbb{F}}
\newcommand{\A}{\mathbb{A}}
\newcommand{\PP}{\mathbb{P}}
\newcommand{\E}{\mathbb{E}}
\newcommand{\eps}{\varepsilon}
\newcommand{\dd}{\,\mathrm{d}}
\begin{document}
\section*{UnsolvedMath ID: 30001306}
\textbf{Source problem code:} catalogue entry
\subsection{Definitions and mathematical statement}
Let \(f\) denote the form in the supplied record, and let \(\pi\) be its attached automorphic representation. For a place \(v\), write \(\pi_v\) for its local component and put
\[
S_{\mathrm{ds}}(\pi)=\{v:\pi_v\text{ is discrete series}\}.
\]
The ground field and the class of forms and representations under consideration are not specified. The notation \(B\) denotes the quaternion algebra implicit in the title, \(f_B\) denotes the associated quaternionic transfer intended by the record, and \(\operatorname{disc}(B)\) denotes its discriminant. The construction and normalization of \(f_B\), the admissible algebras \(B\), and the convention concerning archimedean ramification are not supplied. The brackets \(\langle f,f\rangle\) and \(\langle f_B,f_B\rangle\) denote the pairings appearing in the record; their definitions and normalizations are unspecified.

The supplied conjectural relation is the existence of complex numbers
\[
(c_v)_{v\in S_{\mathrm{ds}}(\pi)},\qquad c_v\in\C,
\]
satisfying
\[
\frac{\langle f,f\rangle}{\langle f_B,f_B\rangle}
\;\sim_p\;
\prod_{v\mid\operatorname{disc}(B)}c_v,
\]
where \(\sim_p\) reproduces the source qualification ``equal up to a \(p\)-adic unit.'' Every index in the product must belong to \(S_{\mathrm{ds}}(\pi)\) for the displayed family to supply its factors, and the denominator must be nonzero for the ratio to be defined. These requirements are not stated explicitly in the record.

This is an unresolved formulation, not a fully specified conjecture. Equality up to a \(p\)-adic unit requires a specified prime, a valued field in which the quantities can be compared, and a means of interpreting the complex period ratio and factors there. If two nonzero quantities \(A,D\) have such an interpretation in a field \(K\) with valuation ring \(\mathcal O_K\), the intended meaning would be \(A/D\in\mathcal O_K^{\times}\). The record supplies neither this comparison framework nor the necessary algebraicity or period normalization. It also does not specify whether one family \((c_v)\) must work simultaneously for every admissible \(B\), or only for one fixed \(B\), nor what local data each factor may depend on. No additional quantifier over \(B\) or hypotheses on \(f\) are asserted here.
\subsection{Short English statement}
The supplied conjecture asks whether the ratio of a form's period pairing to that of its quaternionic transfer factors into complex local constants at the discriminant places, up to a \(p\)-adic unit. The forms, transfers, pairings, admissible quaternion algebras, uniformity of the factors, and \(p\)-adic comparison must be specified before this becomes a precise mathematical assertion.
\subsection{Sources}
\begin{itemize}
\item Local registry: \texttt{data/unsolvedmath/problems.json}, permanent UnsolvedMath ID \(30001306\), exact source problem code \texttt{catalogue entry}, title ``Local Period Factors for Quaternionic Transfers.'' The supplied record belongs to the Oberwolfach Reports Open Problems collection and gives proposed year 2009; no individual proposer is identified.
\item Supplied report citation: \emph{Explicit Methods in Number Theory (2010)}, Oberwolfach Reports, DOI 10.4171/owr/2009/33, pp. 1843--1920. \url{https://doi.org/10.4171/owr/2009/33}. The background describes Oberwolfach Reports as documenting mathematical workshops and their open-problem discussions.
\item Supplied literature source, labeled ``Prasanna Oberwolfach report'': \url{https://ems.press/content/serial-article-files/46233?nt=1}.
\item Supplied literature source, labeled ``Prasanna integral period theorem'': \url{https://doi.org/10.4007/annals.2006.163.901}.
\item Literature assessment dated exactly \texttt{2026-08-21}: the supplied curation labels the problem partially solved and states that the local-factor formulation has substantial results from Harris and later integral work, while the full conjectural relation in the extracted generality is not completely resolved. This is supplied dated curation, not an independently verified literature assessment. The background explicitly requests independent verification of time-sensitive status and cited claims.
\item The registry labels the statement \texttt{corrected\_verified} and describes the local-period-factor conjecture as complete, with representation, discriminant, and product notation normalized. These are supplied source-verification claims; this draft does not independently verify them or resolve the missing hypotheses.
\end{itemize}
\end{document}
